\chapter{2018年上海卷（春）}

\section{填空题}

\begin{problem}
不等式 \(\lvert x\rvert>1\) 的解集为\fillinblank{}.
\end{problem}

\begin{answer}
\((-\infty,-1)\cup(1,+\infty)\)
\end{answer}

\begin{problem}
计算 \(\lim\limits_{n\to\infty}\frac{3n-1}{n+2}=\fillinblank{}\).
\end{problem}

\begin{answer}
\(3\)
\end{answer}

\begin{problem}
设集合 \(A=\{x\mid 0<x<2\}\)，\(B=\{x\mid -1<x<1\}\)，则 \(A\cap B=\fillinblank{}\).
\end{problem}

\begin{answer}
\((0,1)\)
\end{answer}

\begin{problem}
若复数 \(z=1+i\)（\(i\) 是虚数单位），则 \(z+\frac{2}{z}=\fillinblank{}\).
\end{problem}

\begin{answer}
\(2\)
\end{answer}

\begin{problem}
已知 \(\{a_n\}\) 是等差数列，若 \(a_2+a_8=10\)，则 \(a_3+a_5+a_7=\fillinblank{}\).
\end{problem}

\begin{answer}
\(15\)
\end{answer}

\begin{problem}
已知平面上动点 \(P\) 到两个定点 \((1,0)\) 和 \((-1,0)\) 的距离之和等于 \(4\)，则动点 \(P\) 的轨迹方程为\fillinblank{}.
\end{problem}

\begin{answer}
\(\frac{x^2}{4}+\frac{y^2}{3}=1\)
\end{answer}

\begin{problem}
如图，在长方体 \(ABCD-A_1B_1C_1D_1\) 中，\(AB=3\)，\(BC=4\)，\(AA_1=5\)，\(O\) 是 \(A_1C_1\) 的中点，则三棱锥 \(A-A_1OB_1\) 的体积为\fillinblank{}.

\bitmapfigure[max width=0.38\linewidth]{img/2018/shanghai_spring/q7_fig1.png}
\end{problem}

\begin{answer}
\(5\)
\end{answer}

\begin{problem}
某校组队参加辩论赛，从 \(6\) 名学生中选出 \(4\) 人分别担任一，二，三，四辩，若其中学生甲必须参赛且不担任四辩，则不同的安排方法种数为\fillinblank{}（结果用数值表示）.
\end{problem}

\begin{answer}
\(180\)
\end{answer}

\begin{problem}
设 \(a\in\R\)，若 \(\left(x^2+\frac{2}{x}\right)^9\) 与 \(\left(x+\frac{a}{x^2}\right)^9\) 的二项展开式中的常数项相等，则 \(a=\fillinblank{}\).
\end{problem}

\begin{answer}
\(4\)
\end{answer}

\begin{problem}
设 \(m\in\R\)，若 \(z\) 是关于 \(x\) 的方程 \(x^2+mx+m^2-1=0\) 的一个虚根，则 \(\lvert z\rvert\) 的取值范围是\fillinblank{}.
\end{problem}

\begin{answer}
\(\left(\frac{\sqrt3}{3},+\infty\right)\)
\end{answer}

\begin{problem}
设 \(a>0\)，函数 \(f(x)=x+2(1-x)\sin(ax)\)，\(x\in(0,1)\). 若函数 \(y=2x-1\) 与 \(y=f(x)\) 的图像有且仅有两个不同的公共点，则 \(a\) 的取值范围是\fillinblank{}.
\end{problem}

\begin{answer}
\(\left(\frac{11\pi}{6},\frac{19\pi}{6}\right]\)
\end{answer}

\begin{problem}
如图，正方形 \(ABCD\) 的边长为 \(20\) 米，圆 \(O\) 的半径为 \(1\) 米，圆心是正方形的中心，点 \(P\)，\(Q\) 分别在线段 \(AD\)，\(CB\) 上，若线段 \(PQ\) 与圆 \(O\) 有公共点，则称点 \(Q\) 在点 \(P\) 的“盲区”中，已知点 \(P\) 以 \(1.5\) 米/秒的速度从 \(A\) 出发向 \(D\) 移动，同时，点 \(Q\) 以 \(1\) 米/秒的速度从 \(C\) 出发向 \(B\) 移动，则在点 \(P\) 从 \(A\) 移动到 \(D\) 的过程中，点 \(Q\) 在点 \(P\) 的盲区中的时长约为\fillinblank{}秒（精确到 \(0.1\)）.

\bitmapfigure[max width=0.50\linewidth]{img/2018/shanghai_spring/q12_fig1.png}
\end{problem}

\begin{answer}
\(4.4\)
\end{answer}

\section{单选题}

\begin{problem}
下列函数中，为偶函数的是（\quad）

\choices
    {\(y=x^{-2}\)}
    {\(y=x^{\frac13}\)}
    {\(y=x^{-\frac12}\)}
    {\(y=x^3\)}
\end{problem}

\begin{answer}
A
\end{answer}

\begin{problem}
如图，在直三棱柱 \(ABC-A_1B_1C_1\) 的棱所在的直线中，与直线 \(BC_1\) 异面的直线的条数为（\quad）

\FigureLayoutDeclare{img/2018/shanghai_spring/q14_fig1.png}{5.2cm}{14bp 1bp 0bp 0bp}
\bitmapfigure[max width=0.34\linewidth]{img/2018/shanghai_spring/q14_fig1.png}

\choices
    {\(1\)}
    {\(2\)}
    {\(3\)}
    {\(4\)}
\end{problem}

\begin{answer}
C
\end{answer}

\begin{problem}
设 \(S_n\) 为数列 \(\{a_n\}\) 的前 \(n\) 项和，“\(\{a_n\}\) 是递增数列”是“\(\{S_n\}\) 是递增数列”的（\quad）

\choices
    {充分非必要条件}
    {必要非充分条件}
    {充要条件}
    {既非充分又非必要条件}
\end{problem}

\begin{answer}
D
\end{answer}

\begin{problem}
已知 \(A\)，\(B\) 为平面上的两个定点，且 \(\lvert AB\rvert=2\)，该平面上的动线段 \(PQ\) 的端点 \(P\)，\(Q\)，满足 \(\lvert AP\rvert\le5\)，\(\overrightarrow{AP}\cdot\overrightarrow{AB}=6\)，\(\overrightarrow{AQ}=-2\overrightarrow{AP}\)，则动线段 \(PQ\) 所形成图形的面积为（\quad）

\choices
    {\(36\)}
    {\(60\)}
    {\(72\)}
    {\(108\)}
\end{problem}

\begin{answer}
B
\end{answer}

\section{解答题}

\begin{problem}
已知 \(y=\cos x\).

\begin{enumerate}
\item 若 \(f(\alpha)=\frac13\)，且 \(\alpha\in[0,\pi]\)，求 \(f\!\left(\alpha-\frac\pi3\right)\) 的值；
\item 求函数 \(y=f(2x)-2f(x)\) 的最小值.
\end{enumerate}
\end{problem}

\begin{solution}
\begin{enumerate}
\item \(\frac{1+2\sqrt6}{6}\).
\item \(-\frac32\).
\end{enumerate}
\end{solution}

\begin{problem}
已知 \(a\in\R\)，双曲线 \(\Gamma:\frac{x^2}{a^2}-y^2=1\).

\begin{enumerate}
\item 若点 \((2,1)\) 在上，求 \(\Gamma\) 的焦点坐标；
\item 若 \(a=1\)，直线 \(y=kx+1\) 与 \(\Gamma\) 相交于 \(A\)，\(B\) 两点，且线段 \(AB\) 中点的横坐标为 \(1\)，求实数 \(k\) 的值.
\end{enumerate}
\end{problem}

\begin{solution}
\begin{enumerate}
\item \((\sqrt3,0)\)，\((-\sqrt3,0)\).
\item \(\frac{\sqrt5-1}{2}\).
\end{enumerate}
\end{solution}

\begin{problem}
如图，利用“平行于圆锥母线的平面截圆锥面，所得截线是抛物线”的几何原理，某快餐店用两个射灯（射出的光锥为圆锥）在广告牌上投影出其标识，如图1所示，图2是投影出的抛物线的平面图，图3是一个射灯投影的直观图，在图2与图3中，点 \(O\)，\(A\)，\(B\) 在抛物线上，\(OC\) 是抛物线的对称轴，\(OC\perp AB\) 于 \(C\)，\(AB=3\) 米，\(OC=4.5\) 米.

\begin{enumerate}
\item 求抛物线的焦点到准线的距离；
\item 在图3中，已知 \(OC\) 平行于圆锥的母线 \(SD\)，\(AB\)，\(DE\) 是圆锥底面的直径，求圆锥的母线与轴的夹角的大小（精确到 \(0.01^\circ\)）.
\end{enumerate}

\bitmapfigure[max width=0.84\linewidth]{img/2018/shanghai_spring/q19_fig1.png}
\end{problem}

\begin{solution}
\begin{enumerate}
\item \(\frac14\).
\item \(9.59^\circ\).
\end{enumerate}
\end{solution}

\begin{problem}
设 \(a>0\)，函数 \(f(x)=\frac{1}{1+a\cdot2^x}\).

\begin{enumerate}
\item 若 \(a=1\)，求 \(f(x)\) 的反函数 \(f^{-1}(x)\)；
\item 求函数 \(y=f(x)\cdot f(-x)\) 的最大值（用 \(a\) 表示）；
\item 设 \(g(x)=f(x)-f(x-1)\)，若对任意 \(x\in(-\infty,0]\)，\(g(x)\ge g(0)\) 恒成立，求 \(a\) 取值范围.
\end{enumerate}
\end{problem}

\begin{solution}
\begin{enumerate}
\item \(f^{-1}(x)=\log_2\frac{1-x}{x}\)，定义域为 \((0,1)\).
\item \(y_{\max}=\frac{1}{(1+a)^2}\)，当 \(x=0\) 时取得.
\item \((0,\sqrt2]\).
\end{enumerate}
\end{solution}

\begin{problem}
若 \(\{c_n\}\) 是递增数列，数列 \(\{a_n\}\) 满足：对任意 \(n\in\N^*\)，存在 \(m\in\N^*\)，使得 \(\frac{a_m-c_n}{a_m-c_{n+1}}\le0\)，则称 \(\{a_n\}\) 是 \(\{c_n\}\) 的“分隔数列”.

\begin{enumerate}
\item 设 \(c_n=2n\)，\(a_n=n+1\)，证明：数列 \(\{a_n\}\) 是 \(\{c_n\}\) 的分隔数列；
\item 设 \(c_n=n-4\)，\(S_n\) 是 \(\{c_n\}\) 的前 \(n\) 项和，\(d_n=c_{3n-2}\)，判断数列 \(\{S_n\}\) 是否是数列 \(\{d_n\}\) 的分隔数列，并说明理由；
\item 设 \(c_n=aq^{n-1}\)，\(T_n\) 是 \(\{c_n\}\) 的前 \(n\) 项和，若数列 \(\{T_n\}\) 是 \(\{c_n\}\) 的分隔数列，求实数 \(a\)，\(q\) 的取值范围.
\end{enumerate}
\end{problem}

\begin{solution}
\begin{enumerate}
\item 证明略.
\item 不是，反例：\(n=4\) 时，\(m\) 无解.
\item \(\begin{cases}
a>0\\
q\ge2
\end{cases}\).
\end{enumerate}
\end{solution}
